%%
%% This is file `poli-min-exemplo.tex',
%% generated with the docstrip utility.
%%
%% The original source files were:
%%
%% ufrj-poli.dtx  (with options: `poliexemplo')
%% 
%% This is a sample undergraduate project which illustrates the use of the
%% `ufrj' document class with the `ufrj-poli' unit style. It is generated from
%% ufrj-poli.dtx by docstrip.
%% 
%% Copyright (C) 2008-2026 Vicente Helano F. Batista, George O. Ainsworth Jr.,
%% Geraldo Xexéo and Eduardo Mangeli. Released under the GNU General Public
%% License version 3 (see the COPYING.txt file).
%% 
\documentclass[grad]{ufrj}
\usepackage{ufrj-poli}
\addbibresource{exemplo.bib}

\begin{document}
  \title{Título do Projeto de Graduação}
  \foreigntitle{Title of the Undergraduate Project}
  \author{Nome do Autor}{Sobrenome}
  \advisor{Nome do Orientador}{Sobrenome}{D.Sc.}{UFRJ}
  % A folha de aprovacao lista a banca declarada aqui, na ordem: o orientador,
  % que preside a banca, e o primeiro (3.1.2.1.3e).
  \examiner{Nome do Orientador Sobrenome}{D.Sc.}{UFRJ}
  \examiner{Nome do Primeiro Examinador Sobrenome}{D.Sc.}{UFRJ}
  \examiner{Nome do Segundo Examinador Sobrenome}{Ph.D.}{UFRJ}
  % A sigla do departamento: DCC, DES, DEG, DET, DHIMA, DEE, DEM, DMM, DNC,
  % DENO, DEI, ECI ou DEL.
  \department{ECI}
  \date{01}{2026}
  \approvaldate{15 de janeiro de 2026}
  \keyword{primeira palavra-chave}
  \keyword{segunda palavra-chave}
  \foreignkeyword{first keyword}
  \foreignkeyword{second keyword}

  % Ordem dos elementos pre-textuais (Manual 3.1.2), que a classe confere:
  % \maketitle, \frontmatter, o resumo em portugues, o abstract, as listas e
  % \tableofcontents, sempre o ultimo.
  \maketitle
  \frontmatter
  \begin{abstract}
    O resumo em português, em parágrafo único, com até 500 palavras
    (3.1.2.1.4).
  \end{abstract}
  \begin{foreignabstract}
    The abstract in English, the version of the abstract in the language of
    international dissemination (3.1.2.1.5).
  \end{foreignabstract}
  \listoffigures
  \tableofcontents

  \mainmatter
  \chapter{Introdução}
  O texto do trabalho começa aqui, e as citações seguem a classe
  \cite{m-4211}.

  \begin{figure}[ht]
    \caption{Uma figura de exemplo}
    \centering
    \rule{6cm}{3cm}
    \source{Elaborado pelo autor (2026).}
  \end{figure}

  \chapter{Conclusão}
  A conclusão do trabalho.

  \backmatter
  \printbibliography
\end{document}
%% 
%%
%% End of file `poli-min-exemplo.tex'.
